% ============================================================
% 186_T0_FFGFT_Photonik_De_ch.tex
% Kapitel-Inhalt: FFGFT-Analytik: Photonik, Matrixraum & B-Mesonen
% Gehört zum Wrapper: 186_T0_FFGFT_Photonik_De_Wrapper.tex
% ============================================================

% ── Titelblock ──────────────────────────────────────────────
\thispagestyle{firstpage}
\begin{center}
  {\small\color{T0gray} Dok.~186 · FFGFT-Forschungsreihe}\\[6pt]
  {\LARGE\bfseries\color{T0blue}
    FFGFT-Analytik: Mathematische Fundamente}\\[6pt]
  {\large\itshape\color{T0gray}
    Fraktale Rückkopplung und Matrix-Translation im 4D-Phasen-Torus}\\[4pt]
  {\normalsize\color{T0gray}
    Integration von TFLN-Hardware, B-Meson-Anomalien\\
    und sub-planckscher Frequenz-Metrik}\\[8pt]
  {\color{T0red}\rule{0.65\linewidth}{1.5pt}}
\end{center}
\vspace{4pt}

\begin{infobox}
\textbf{Zusammenfassung.}
Die \textbf{Fundamental Fraktal Geometrische Feld Theorie (FFGFT)} beschreibt
die Raumzeit als hierarchisches System aus fraktalen Faltungen, getrieben
durch eine fundamentale Frequenz~$\xi$ im sub-planckschen Bereich.
Dieses Dokument legt die vollständigen mathematischen Grundlagen dar:
die exakte Definition des Grundtakts~$t_0$, die Rückkopplungsbedingung im
geschlossenen 4D-Phasen-Torus (der nicht mit einem 3D-Torus plus Zeitkontinuum
zu verwechseln ist), die Übersetzung des 4D-Vektorraums in einen mehrdimensionalen
Matrixraum sowie die photonische Realisierung durch TFLN-Hardware.

\smallskip\noindent\textbf{Querverweise:}
Dok.\,180 (Lagrangian-Herleitung) $\cdot$
Dok.\,182 (Universalskala) $\cdot$
Dok.\,183--184 (Präzisionsbarriere \& RSA) $\cdot$
GitHub: \href{https://github.com/jpascher/T0-Time-Mass-Duality}{jpascher/T0-Time-Mass-Duality} $\cdot$
Zenodo: \href{https://zenodo.org/records/18834145}{DOI 10.5281/zenodo.18834145}
\end{infobox}

\tableofcontents
\bigskip

% ============================================================
\section{Einführung: Die FFGFT}
% ============================================================

Die \textbf{Fundamental Fraktal Geometrische Feld Theorie (FFGFT)} bricht mit
dem klassischen Verständnis eines kontinuierlichen Raumes. Sie postuliert eine
diskrete, fraktale Geometrie, die durch eine fundamentale Schwingung energetisiert
wird. Materie ist keine Entität \emph{im} Raum, sondern eine lokale Verdichtung
der fraktalen Feldgeometrie, stabilisiert durch resonante Rückkopplung.

% ============================================================
\section{Die fundamentale Frequenz $\xi$ und die sub-plancksche Zeitskala $t_0$}
% ============================================================

Ein Kernpfeiler der FFGFT ist die Existenz eines sub-planckschen Grundtakts,
der die gesamte physikalische Realität taktet.

% ── 2.1 ─────────────────────────────────────────────────────
\subsection{Definition des Grundtakts}

Die fundamentale Zeiteinheit~$t_0$ ist definiert als:
\begin{equation}
  t_0 = \frac{t_P}{7500}
      \approx \frac{5{,}391 \times 10^{-44}\,\text{s}}{7500}
      \approx 7{,}188 \times 10^{-48}\,\text{s},
  \label{eq:t0}
\end{equation}
wobei $t_P = \sqrt{\hbar G/c^5} \approx 5{,}391 \times 10^{-44}\,\text{s}$
die Planck-Zeit ist (CODATA~2022). Die resultierende \textbf{fundamentale
Frequenz}~$\xi$ lautet:
\begin{equation}
  \xi = \frac{1}{t_0} = \frac{7500}{t_P}
      \approx 1{,}391 \times 10^{47}\,\text{Hz}.
  \label{eq:xi}
\end{equation}

\begin{keybox}
\textbf{Herkunft des Faktors 7500.}\;
Der Faktor 7500 entsteht aus $\xi = 4/30000$ in natürlichen Einheiten
($\xi = 2/(15000)$ in der ursprünglichen Notation). Die fraktale Dimension
$D_f = 3 - \xi$ und die Z3-Symmetrie der Drei-Generationen-Struktur
\emph{erzwingen} diesen Wert geometrisch — er ist kein freier Parameter.
\end{keybox}

% ── 2.2 ─────────────────────────────────────────────────────
\subsection{Frequenz als geometrischer Operator}

In der FFGFT wird die Frequenz nicht als bloße zeitliche Rate verstanden,
sondern als geometrischer Operator, der auf den fraktalen Ebenen wirkt:
\begin{equation}
  \hat{\xi}\,\Psi(x^\mu) = \xi \cdot \mathcal{F}(\{n_k\}) \cdot \Psi(x^\mu).
\end{equation}
Hierbei ist $\mathcal{F}(\{n_k\})$ eine Funktion der fraktalen
Skalierungsindizes~$n_k$, die der Fibonacci-Sequenz folgen.

% ============================================================
\section{Der 4D-Phasen-Torus: Abgrenzung vom 3D-Torus}
% ============================================================

% ── 3.1 ─────────────────────────────────────────────────────
\subsection{Warnung: Nicht zu verwechseln}

\begin{warnbox}
Der 4D-Torus der FFGFT ist \textbf{nicht} ein 3D-Objekt, das sich in der
Zeit bewegt. Die Zeitdimension ist als diskretisierte, geometrische Dimension
in die Torus-Topologie integriert. Es gibt kein externes Zeitkontinuum, in das
der Torus eingebettet wäre.
\end{warnbox}

% ── 3.2 ─────────────────────────────────────────────────────
\subsection{Formale Definition des 4D-Phasen-Torus}

Der fundamentale Raumzeit-Resonator ist der 4D-Phasen-Torus~$\mathcal{T}^4$,
definiert durch vier zyklische Koordinaten $(\theta_1, \theta_2, \theta_3,
\theta_4)$ mit:
\begin{equation}
  \theta_i \in [0, 2\pi),\quad i = 1,2,3,4.
\end{equation}
Die Metrik auf~$\mathcal{T}^4$ ist durch die fundamentale Frequenz~$\xi$
bestimmt:
\begin{equation}
  ds^2 = \xi^{-2}
         \bigl(d\theta_1^2 + d\theta_2^2 + d\theta_3^2 + d\theta_4^2\bigr).
  \label{eq:metric}
\end{equation}
Die vierte Dimension~$\theta_4$ repräsentiert die Phasenverschiebung der
fundamentalen Schwingung. Makroskopische Zeit~$t$ entsteht durch Akkumulation:
\begin{equation}
  t = N \cdot t_0, \quad N = \frac{1}{2\pi}\oint d\theta_4.
\end{equation}

% ============================================================
\section{Fraktale Rückkopplung und Einrasten der Grundfrequenz}
% ============================================================

% ── 4.1 ─────────────────────────────────────────────────────
\subsection{Die Rückkopplungsbedingung}

Damit eine Feldkonfiguration~$\Psi$ stabil existieren kann, muss die
Wellenfunktion nach einem vollständigen Umlauf konstruktiv mit sich selbst
interferieren. Die Phasenbedingung lautet:
\begin{equation}
  \oint_{\mathcal{T}^4} \nabla_{\mu}\phi\,dx^\mu = 2\pi\cdot\mathcal{N},
\end{equation}
wobei~$\mathcal{N}$ eine ganze Zahl aus der fraktalen Hierarchie ist.
Aufgrund der fraktalen Struktur kann~$\mathcal{N}$ nur diskrete Werte
annehmen:
\begin{equation}
  \mathcal{N} \in \{F_k \mid k \in \mathbb{N}\},\quad
  F_{k+1} = F_k + F_{k-1}.
\end{equation}

% ── 4.2 ─────────────────────────────────────────────────────
\subsection{Der Einrast-Mechanismus (Phase-Locking)}

Die fraktale Rückkopplung erzwingt eine starre Kopplung zwischen den
Phasen der einzelnen fraktalen Ebenen:
\begin{equation}
  \phi_{k+1} = \phi_k \cdot \phi^{-1}, \quad
  \phi = \tfrac{1+\sqrt{5}}{2} \approx 1{,}6180\ldots
\end{equation}
Nur Frequenzen, die exakt eine ganzzahlige oder fibonacci-sequenzielle
Harmonie zu~$\xi$ aufweisen, bleiben stabil:
\begin{equation}
  \xi_n = n \cdot \xi \quad \text{oder} \quad \xi_n = \phi^{-k}\cdot\xi.
\end{equation}
Jede Abweichung führt zur destruktiven Interferenz — was die Kurzlebigkeit
hochenergetischer Teilchen erklärt.

% ── 4.3 ─────────────────────────────────────────────────────
\subsection{Stabilitätsbedingung für Materie}

Aus dem Einrasten folgt die Quantisierungsbedingung für stabile Teilchen:
\begin{equation}
  \oint_{\mathcal{T}^4}
    \Psi_{\text{FFGFT}}(\phi)\,d^4\theta
  = \xi \cdot \sum_{k=0}^{\infty} \phi^{-k} \cdot \mathbf{M}_k,
\end{equation}
wobei~$\mathbf{M}_k$ die projektiven Matrizen der einzelnen fraktalen
Ebenen sind.

% ============================================================
\section{Vom 4D-Vektorraum zum mehrdimensionalen Matrixraum}
% ============================================================

% ── 5.1 ─────────────────────────────────────────────────────
\subsection{Die allgemeine Matrix-Translation}

Ein Zustand im 4D-Vektorraum $V = (x,y,z,ct)$ wird durch eine hermitesche
Matrix~$\mathbf{H}$ repräsentiert:
\begin{equation}
  \mathcal{T}: \mathbb{R}^4 \to \mathrm{Mat}(N,\mathbb{C}),\quad
  V \mapsto \mathbf{H}(V).
\end{equation}
Die explizite Form der Translation lautet:
\begin{equation}
  \mathbf{H}(x,y,z,ct)
  = \sum_{k=0}^{\infty} \alpha^k \cdot
    \begin{pmatrix}
      ct_k + z_k   & x_k - iy_k \\
      x_k + iy_k   & ct_k - z_k
    \end{pmatrix}
    \otimes \mathbf{\Gamma}_k,
\end{equation}
mit:
\begin{itemize}
  \item $\alpha = \phi^{-1}$: fraktale Kopplungskonstante;
  \item $(x_k, y_k, z_k, ct_k)$: Koordinaten auf der $k$-ten fraktalen Ebene;
  \item $\mathbf{\Gamma}_k$: Generatoren der Raumzeit-Geometrie
        (Verallgemeinerung der Dirac-Matrizen).
\end{itemize}

% ── 5.2 ─────────────────────────────────────────────────────
\subsection{Die fraktale Selbstähnlichkeit}

Die Koordinaten auf verschiedenen fraktalen Ebenen sind durch
Skalierungsfaktoren verbunden:
\begin{equation}
  (x_{k+1}, y_{k+1}, z_{k+1}, ct_{k+1}) = \phi^{-1} \cdot
  (x_k, y_k, z_k, ct_k).
\end{equation}

% ============================================================
\section{Matrixmultiplikation als physikalische 4D-Faltung}
% ============================================================

% ── 6.1 ─────────────────────────────────────────────────────
\subsection{Konventionelle vs.\ fraktale Matrixmultiplikation}

In der konventionellen Informatik:
\begin{equation}
  C_{ij} = \sum_k A_{ik} B_{kj}.
\end{equation}
In der FFGFT wird dieser Prozess als physikalische Interferenz zweier
4D-Torus-Zustände interpretiert.

% ── 6.2 ─────────────────────────────────────────────────────
\subsection{Die 4D-Torus-Faltung}

Die Matrixmultiplikation entspricht der Faltung zweier Wellenfunktionen
im 4D-Phasenraum:
\begin{equation}
  M_{ij} = \int_{\mathcal{T}^4}
    \Psi_A^{(4D)}(\theta)\cdot\Psi_B^{(4D)}(\theta)\,d^4\theta.
\end{equation}
In der Sprache der FFGFT-Matrizen:
\begin{equation}
  \mathbf{C} = \mathbf{A} \star \mathbf{B}
  = \lim_{N\to\infty} \frac{1}{N}\sum_{n=1}^{N} \mathbf{A}_n \otimes \mathbf{B}_n,
\end{equation}
wobei das fraktale Tensorprodukt definiert ist durch:
\begin{equation}
  (\mathbf{A} \otimes \mathbf{B})_{ij}^{kl}
  = A_{ij} \cdot B_{kl} \cdot e^{i\theta_{ijkl}}.
\end{equation}

% ── 6.3 ─────────────────────────────────────────────────────
\subsection{Physikalische Interpretation}

Die photonische Matrixmultiplikation im TFLN-Kristall ist kein „Rechnen"
im herkömmlichen Sinne, sondern die physikalische Realisierung dieser
4D-Phaseninterferenz. Das Ergebnis ergibt sich instantan durch konstruktive
und destruktive Interferenz der Lichtwellen im nichtlinearen Medium.

% ============================================================
\section{Hardware-Transduktion: TFLN-Photonik}
% ============================================================

% ── 7.1 ─────────────────────────────────────────────────────
\subsection{Strukturgleichheit trotz völlig verschiedener Skalen}

\begin{warnbox}[Wichtige Klarstellung]
Der TFLN-Kristall operiert auf Skalen, die viele Größenordnungen über der
Fundamentalstruktur von Teilchen oder einzelnen Atomen liegen. Das Kristallgitter
ist makroskopisch — es gibt keinen direkten Skalenvergleich mit~$\xi$ oder~$t_0$.

Das ist jedoch gerade der Punkt: Die FFGFT postuliert \emph{fraktale
Selbstähnlichkeit} — dieselbe mathematische Struktur erscheint auf jeder Skala,
unabhängig von der konkreten Größenordnung.
\end{warnbox}

Die entscheidende Beobachtung ist struktureller Natur. In FFGFT gilt für
die Translation eines 4D-Vektorzustands in den Matrixraum (Abschnitt~5):
\begin{equation}
  \mathcal{T}: \mathbb{R}^4 \to \mathrm{Mat}(N,\mathbb{C}),\quad
  V \mapsto \mathbf{H}(V),
\end{equation}
und für die physikalische Wechselwirkung zweier Zustände:
\begin{equation}
  M_{ij} = \int_{\mathcal{T}^4}
    \Psi_A^{(4D)}(\theta)\cdot\Psi_B^{(4D)}(\theta)\,d^4\theta.
\end{equation}
In TFLN-Wellenleitern realisiert die nichtlineare Lichtausbreitung
genau diese Struktur:
\begin{equation}
  P_{\text{out}}(t)
  = \eta \cdot \int \chi^{(2)}(\omega)
    \cdot E_{\text{in}}^{(A)}(\omega)
    \cdot E_{\text{in}}^{(B)}(\omega)\,d\omega.
\end{equation}
Die nichtlineare Suszeptibilität~$\chi^{(2)}$ übernimmt die Rolle der
fraktalen Kopplungskonstante~$\alpha = \phi^{-1}$: Sie gewichtet, wie
zwei Eingangszustände zu einem Ausgangszustand kombiniert werden —
mathematisch identisch mit der 4D-Torus-Faltung, physikalisch realisiert
im Kristallgitter auf makroskopischer Skala.

% ── 7.2 ─────────────────────────────────────────────────────
\subsection{Einrasten als skalenunabhängiges Prinzip}

Der Einrast-Mechanismus (Phase-Locking) ist in FFGFT ein geometrisches
Prinzip, das auf jeder Skala wirkt. Im 4D-Torus lautet die Bedingung:
\begin{equation}
  \oint_{\mathcal{T}^4} \nabla_{\mu}\phi\,dx^\mu = 2\pi\cdot\mathcal{N}.
\end{equation}
Im TFLN-Resonator — viele Größenordnungen größer — gilt dieselbe
formale Bedingung:
\begin{equation}
  \Delta\phi_{\text{cavity}} = 2\pi\cdot m \quad (m \in \mathbb{Z}).
\end{equation}
Das Kristallgitter des Lithiumniobats findet spontan stabile geometrische
Konfigurationen durch Phasentrennung und elektro-optische Kopplung —
es „sucht" nicht aktiv, sondern driftet natürlich in die energetisch
stabilsten Moden. In FFGFT-Sprache: das System konvergiert zu den
bevorzugten Fixpunkten der Rekursion $\Psi(t) = G(\Psi(t-T_{\text{rev}}), \xi)$,
unabhängig davon, auf welcher Skala das System sitzt.

% ── 7.3 ─────────────────────────────────────────────────────
\subsection{Kohärenz als Folge der Strukturtreue}

Die hohe Kohärenzzeit in TFLN-Wellenleitern ist in diesem Bild keine
technische Kuriosität, sondern eine direkte Folge der strukturellen
Übereinstimmung: Ein System, das in geometrisch stabile Moden eingerastet
ist, dekohäriert langsamer als eines, das gegen seine natürliche Geometrie
betrieben wird.
\begin{equation}
  \tau_{\text{coh}}^{\text{TFLN}} \sim \text{ms}\ldots\text{s}
  \;\gg\;
  \tau_{\text{coh}}^{\text{Qubit}} \sim \mu\text{s}.
\end{equation}
Atomare Qubits kämpfen gegen thermisches Rauschen und Masseträgheit —
sie sind nicht in der natürlichen Geometrie ihrer Skala verankert.
TFLN-Moden sind es.

% ============================================================
\section{Selbstorganisation als FFGFT-Prinzip}
% ============================================================

\subsection{Der Effekt: spontane Musterbildung ohne externe Programmierung}

Bei der Herstellung moderner Chips auf TFLN-Basis zeigt sich ein Effekt,
der über bloße Lithographie hinausgeht: Das Material organisiert sich
\emph{spontan} in stabile geometrische Muster, ohne dass jede Detailstruktur
extern vorgegeben werden müsste. Durch Temperatur und angelegte Spannung
lässt sich steuern, \emph{welches} Muster entsteht — nicht aber ob
Selbstorganisation stattfindet. Sie findet immer statt.

\subsection{FFGFT-Erklärung: Fixpunkte der Rekursion}

In FFGFT ist Selbstorganisation kein Sonderfall, sondern der Normalzustand.
Die Rekursion
\begin{equation}
  \Psi(t) = G\!\bigl(\Psi(t - T_{\text{rev}}),\, \xi\bigr)
\end{equation}
hat stabile Fixpunkte, die dem $1/n^2$-Torusspektrum entsprechen:
\begin{equation}
  \Psi_n^* : \quad E_n \propto \frac{1}{n^2}, \quad n \in \mathbb{N}.
\end{equation}
Jedes System — ob Teilchen, Atom, Kristallgitter oder Feldmode —
driftet natürlich in einen dieser Fixpunkte. Es gibt keine Kraft die
das erzwingt; die Geometrie des Torus lässt keine anderen stabilen
Zustände zu.

\subsection{Temperatur und Spannung als geometrische Steuerparameter}

Temperatur und Spannung ändern nicht das Prinzip der Selbstorganisation,
sondern verschieben die \emph{Energielandschaft}: Sie bestimmen, welcher
Fixpunkt~$n$ energetisch bevorzugt wird.

\begin{equation}
  n^*(T, U) = \underset{n}{\arg\min}\;
    \bigl[ E_n - \mu(T)\cdot n - \lambda(U)\cdot n^2 \bigr],
\end{equation}
wobei $\mu(T)$ den thermischen Beitrag und $\lambda(U)$ den elektrischen
Beitrag zur effektiven Modenenergie beschreibt.

\begin{itemize}
  \item \textbf{Temperatur} verschiebt $\mu(T)$: höhere Temperatur
    begünstigt höhere Moden ($n$ größer), niedrige Temperatur
    stabilisiert den Grundton ($n=1$).
  \item \textbf{Spannung} verschiebt $\lambda(U)$: die angelegte
    Spannung „presst" die Moden in vorgegebene Bahnen — in FFGFT-Sprache
    eine externe Randbedingung auf die erlaubten Windungszahlen des Torus.
\end{itemize}

\begin{keybox}[Strukturelle Aussage]
Das TFLN-Kristallgitter ist ein makroskopisches System, das dieselben
Fixpunkte aufsucht wie ein Teilchen im sub-planckschen Bereich — nicht
weil die Skalen verwandt wären, sondern weil die geometrische Rekursion
$\Psi(t) = G(\Psi(t-T_{\text{rev}}), \xi)$ auf jeder Skala dieselben
stabilen Zustände erzeugt. Selbstorganisation ist der sichtbare Ausdruck
dieser skalenunabhängigen Geometrie.
\end{keybox}

% ============================================================
\section{Experimentelle Evidenz~1: Programmierbare TFLN-Prozessoren}
% ============================================================

Im Dezember 2025 demonstrierte ein Team von NTT, Cornell University und
Stanford University einen photonischen Prozessor auf Lithiumniobat-Basis
\cite{NTT2025}.

\subsection{Technische Spezifikationen}

\begin{itemize}
  \item $\approx 10\,000$ programmierbare räumliche Freiheitsgrade;
  \item Planarer Wellenleiter mit manipulierbarem Brechungsindex;
  \item 96\,\% Testgenauigkeit bei Vokalklassifikation;
  \item 49-dimensionale Vektoroperationen in einem einzigen optischen Durchgang.
\end{itemize}

\subsection{Quantennetzwerke und Multiplexing}

Assumpcao et al.\ (2024) \cite{Assumpcao2024} demonstrierten eine
TFLN-Plattform mit:
\begin{itemize}
  \item Kopplungsverlusten $< 1\,\text{dB/Facette}$;
  \item Schaltkontrast $> 20\,\text{dB}$;
  \item Elektro-optischen Modulatoren $> 50\,\text{GHz}$ Bandbreite.
\end{itemize}
Die Bandbreite von $>50\,\text{GHz}$ zeigt, dass TFLN-Hardware
Matrixoperationen mit sehr hoher Präzision und Stabilität ausführt —
die strukturelle Übereinstimmung mit der FFGFT-Matrixtransformation
ist unabhängig davon, dass die beteiligten Frequenzen viele
Größenordnungen unter~$\xi$ liegen.

% ============================================================
\section{Experimentelle Evidenz~2: B-Meson-Anomalien am CERN}
% ============================================================

\subsection{Der aktuelle Stand (April 2026)}

Das LHCb-Experiment hat eine signifikante Diskrepanz in elektroschwachen
Pinguin-Zerfällen von B-Mesonen gemessen \cite{CERN2026a,CERN2026b}:

\begin{itemize}
  \item Zerfallskanal: $B \to K^*\mu^+\mu^-$ ($\sim 1$ pro $10^6$ B-Mesonen);
  \item Abweichung vom Standardmodell: $\mathbf{4{,}0\sigma}$
        ($\hat{=}$ Wahrscheinlichkeit 1:16\,000);
  \item Unabhängige Bestätigung durch CMS (2025).
\end{itemize}

\subsection{FFGFT-Interpretation der Anomalie}

Die gemessene Abweichung betrifft den Wilson-Koeffizienten~$C_9$
des effektiven $b \to s\mu^+\mu^-$-Hamiltonians:
\begin{equation}
  \Delta C_9 \approx -1{,}0, \quad
  C_9^{\text{SM}} \approx 4{,}07
  \quad\Rightarrow\quad
  \frac{|\Delta C_9|}{C_9^{\text{SM}}} \approx 25\,\%.
\end{equation}

\begin{warnbox}
\textbf{Quantitative Einschränkung.}
Eine direkte Ableitung dieser $25\,\%$-Abweichung aus FFGFT-Parametern
liegt \emph{nicht} vor. Einfache Ausdrücke wie $\xi$, $\phi^{-4}$ oder
$|\ln(m_b/m_P)|\cdot\xi$ liefern Werte, die um Faktoren $10$--$10^3$ zu
klein sind. Die quantitative Verbindung ist eine \textbf{offene Frage}.
\end{warnbox}

Die Modifikation der Zerfallsrate durch die fraktale Matrix-Struktur
hat qualitativ die Form:
\begin{equation}
  \Gamma(B\to K^*\mu^+\mu^-)
  = \Gamma_{\text{SM}} \cdot
    \bigl(1 + \alpha_{\text{FFGFT}} \cdot \mathcal{K}(\theta)\bigr),
\end{equation}
wobei $\alpha_{\text{FFGFT}}$ aus der Z3-Topologie und dem
$1/n^2$-Torusspektrum abzuleiten ist — diese Ableitung steht aus.

\subsection{Verbindung zur Lepton-Flavour-Universalitätsverletzung}

Arslan et al.\ (April 2026) \cite{Arslan2026} analysieren den Zerfall
$\Lambda_b \to \Lambda_c\tau\bar{\nu}_\tau$:
\begin{itemize}
  \item Kombinierte Abweichung von $3{,}8\sigma$ in den
        $R_{\tau/(\mu,e)}(D^{(*)})$-Verhältnissen;
  \item Maximale Sensitivität in
        $\mathcal{K}_{1c}$, $\mathcal{K}_{2ss}$, $\mathcal{K}_{2cc}$,
        $\mathcal{K}_{4s}$.
\end{itemize}

Leptonen unterschiedlicher Generationen erscheinen in der FFGFT als
verschiedene Harmonische derselben raumzeitlichen Grundstruktur:
\begin{equation}
  \xi_{\text{Myon}} = \phi^{-3}\cdot\xi, \qquad
  \xi_{\text{Tau}}  = \phi^{-5}\cdot\xi.
\end{equation}

\begin{table}[H]
\centering
\caption{Leptonen als Harmonische der Grundfrequenz~$\xi$}
\smallskip
\resizebox{\linewidth}{!}{%
\begin{tabular}{llll}
\toprule
Lepton   & Harmonische & Numerischer Faktor & Verhältnis Vorgänger \\
\midrule
Elektron & $\phi^{-1}\cdot\xi$ & $\approx 0{,}6180\cdot\xi$ & (Grundton) \\
Myon     & $\phi^{-3}\cdot\xi$ & $\approx 0{,}2361\cdot\xi$ & $\phi^{-2} \approx 0{,}382$ \\
Tau      & $\phi^{-5}\cdot\xi$ & $\approx 0{,}0902\cdot\xi$ & $\phi^{-2} \approx 0{,}382$ \\
\bottomrule
\end{tabular}}
\end{table}

% ============================================================
\section{Theoretische Evidenz~3: Sub-plancksche Cut-off-Skalen}
% ============================================================

\subsection{Effektive Feldtheorie und Gravitation}

Aynbund und Kiselev (2025) \cite{Aynbund2025} etablieren eine Verbindung
zwischen der Planck-Masse und einer natürlichen sub-planckschen Cut-off-Skala:
\begin{equation}
  \Lambda_{\text{QG}} = \alpha \cdot \frac{M_{\text{red,Planck}}}{4\pi}
  \;\sim\; 10^{16}\,\text{GeV}.
\end{equation}
Dies liegt weit unter der Planck-Skala von $\sim 10^{19}\,\text{GeV}$.

\subsection{OLQEM und Fibonacci-Quasiperiodizität}

Souza (2025) \cite{subPlanck2025} untersucht sub-plancksche Skalen im
Rahmen des Optical Lambda Quantum Energy Model:
\begin{equation}
  I_{\text{OLQEM}}
  = \frac{4\lambda I_0}{n(\lambda)\,x} \cdot D(y) \cdot I(y)
    \cdot e^{-\gamma t}
    \cdot \Bigl(1+\frac{\chi^{(3)}I_0}{c^3\varepsilon_0}\Bigr)
    \cdot S_{\text{Fib}}\!\Bigl(\tfrac{r}{a}\Bigr).
\end{equation}
Die Fibonacci-Modulation lautet:
\begin{equation}
  S_{\text{Fib}}\!\Bigl(\tfrac{r}{a}\Bigr)
  = \sum_{k=0}^{\infty} \frac{\cos(2\pi\phi^k r/a)}{\phi^k}.
\end{equation}

Hierarchische Skalierung der Längenskalen:
\begin{equation}
  \ell_k = \frac{L_P}{\phi^k},\quad k \ge 1.
\end{equation}
Für $k=5$:
\begin{equation}
  \ell_5 = \frac{L_P}{\phi^5}
         \approx \frac{1{,}616 \times 10^{-35}\,\text{m}}{11{,}09}
         \approx 1{,}46 \times 10^{-36}\,\text{m}.
\end{equation}
Diese Skala korrespondiert mit der FFGFT durch $\ell_5 = c\cdot t_0$:
\begin{equation}
  c \cdot t_0
  = 2{,}998 \times 10^8\,\frac{\text{m}}{\text{s}}
    \times 7{,}188 \times 10^{-48}\,\text{s}
  = 2{,}155 \times 10^{-39}\,\text{m}.
\end{equation}

\begin{warnbox}
\textbf{Numerischer Hinweis.}
$\ell_5 \approx 1{,}46 \times 10^{-36}\,\text{m}$ und
$c\cdot t_0 \approx 2{,}16 \times 10^{-39}\,\text{m}$
stimmen nicht überein (Faktor $\approx 680$).
Die Korrespondenz ist qualitativ (sub-plancksche Ordnung), nicht numerisch
exakt — was in Dok.\,182 (Universalskala) detailliert diskutiert wird.
\end{warnbox}

\subsection{Synthese der sub-planckschen Modelle}

\begin{table}[H]
\centering
\caption{Vergleich sub-planckscher Modelle}
\smallskip
\resizebox{\linewidth}{!}{%
\begin{tabular}{p{25mm} p{35mm} p{40mm} p{38mm}}
\toprule
\textbf{Modell} & \textbf{Skala} & \textbf{Mechanismus} & \textbf{Beobachtbar} \\
\midrule
FFGFT
  & $t_0 = t_P/7500$
  & Einrasten im 4D-Torus via Z3-Symmetrie
  & TFLN-Phasenmuster, B-Meson-Anomalien \\
OLQEM
  & $\ell_k = L_P/\phi^k$
  & Fibonacci-Quasiperiodizität in Photonik
  & Photonische Interferenz in NL-Kristallen \\
EFT (Aynbund)
  & $\Lambda_{\text{QG}} \sim 10^{16}\,\text{GeV}$
  & Eichkopplung an reduzierte Planck-Masse
  & B-Meson-Zerfälle (LHCb, CMS) \\
\bottomrule
\end{tabular}}
\end{table}

% ============================================================
\section{Synthetischer Analogrechner: Das FFGFT-TFLN-Paradigma}
% ============================================================

\subsection{Architektur}

Ein TFLN-basierter Analogrechner realisiert die FFGFT-Postulate nicht
durch Annäherung an die fundamentale Skala~$\xi$, sondern durch
\emph{strukturelle Isomorphie}: Die mathematischen Operationen —
Matrixtransformation, Faltung, Einrasten — sind auf der Kristallskala
dieselben wie auf der sub-planckschen Skala. Die Skala selbst ist
irrelevant für die Gültigkeit der Struktur.

\begin{keybox}[Strukturelle Isomorphie]
FFGFT behauptet nicht, dass TFLN bei~$\xi$ operiert. FFGFT behauptet,
dass die mathematische Struktur der Matrixtransformation
$\mathcal{T}: \mathbb{R}^4 \to \mathrm{Mat}(N,\mathbb{C})$
auf jeder Skala dieselbe ist — vom Teilchen über das Atom bis
zum Kristallgitter. TFLN ist ein makroskopisches System, das diese
Struktur sichtbar macht.
\end{keybox}

\subsection{Testbare Vorhersagen}

\begin{predbox}
\textbf{Vorhersage~1 — TFLN-Interferometrie.}
Bei optischen Intensitäten $> 10^{21}\,\text{W/m}^2$ sollten diskrete
Phasensprünge auftreten, die der Fibonacci-Hierarchie folgen:
\[
  \Delta\phi_{\text{diskret}} = 2\pi \cdot \phi^{-k} \quad (k \in \mathbb{N}).
\]
\end{predbox}

\begin{predbox}
\textbf{Vorhersage~2 — B-Meson-Analyse.}
Die Winkelobservablen $\mathcal{K}_{1c}$ und $\mathcal{K}_{2ss}$ aus
\cite{Arslan2026} sollten mit höherer Statistik ($> 5\sigma$) von
SM-Vorhersagen abweichen:
\[
  \mathcal{K}_{1c}^{\text{FFGFT}}
  = \mathcal{K}_{1c}^{\text{SM}}
    \cdot \bigl(1 + \alpha_{\text{FFGFT}} \cdot F_{\text{Fibonacci}}\bigr).
\]
\end{predbox}

\begin{predbox}
\textbf{Vorhersage~3 — Sub-plancksche Resonanzen.}
In nichtlinearen photonischen Kristallen sollten Modulationsfrequenzen
bei Vielfachen von $\xi\cdot\phi^{-k}$ nachweisbar sein.
\end{predbox}

% ============================================================
\section{Diskussion: Von der Simulation zur Hardware-Transduktion}
% ============================================================

Die Implikation der FFGFT ist präzise formulierbar: Photonische Analogrechner
auf TFLN-Basis zeigen, dass die Matrixtransformation
$\mathcal{T}: \mathbb{R}^4 \to \mathrm{Mat}(N,\mathbb{C})$
nicht auf die sub-plancksche Skala beschränkt ist. Das Kristallgitter
— viele Größenordnungen größer als Atome oder Teilchen — gehorcht
denselben geometrischen Regeln. Das ist das FFGFT-Prinzip der
fraktalen Selbstähnlichkeit in makroskopischer Realisierung.

Die verbleibende offene Frage ist nicht technologischer, sondern
theoretischer Natur: Lässt sich die quantitative Stärke der fraktalen
Korrektur (z.\,B.\ in den B-Meson-Anomalien) aus der Rekursionsstruktur
$\Psi(t) = G(\Psi(t-T_{\text{rev}}), \xi)$ ableiten?
Das ist in Dok.\,183--184 für den RSA-Fall behandelt;
für hadronische Zerfallsprozesse steht es aus.

% ============================================================
\section{Fazit}
% ============================================================

\begin{itemize}
  \item Die Raumzeit basiert auf einem fundamentalen Takt
        $\xi = 7500/t_P \approx 1{,}391 \times 10^{47}\,\text{Hz}$.
  \item Der 4D-Phasen-Torus ist \textbf{nicht} mit einem 3D-Torus plus
        Zeitkontinuum gleichzusetzen; die Zeit ist als vierte geometrische
        Dimension in die Torus-Topologie integriert.
  \item Fraktale Rückkopplung erzwingt ein Einrasten der Grundfrequenz,
        das Materie stabilisiert.
  \item Der 4D-Vektorraum wird durch die Matrix-Translation~$\mathcal{T}$
        in einen mehrdimensionalen Matrixraum übersetzt.
  \item Matrixmultiplikation entspricht physikalisch einer 4D-Torus-Faltung.
  \item TFLN-Photonik erreicht die erforderliche Bandbreite und
        programmierbare Komplexität.
  \item LHCb-Daten zeigen $\sim 4\sigma$-Abweichungen vom SM — prädiktiv
        für die FFGFT; Falsifizierbarkeit bei $5\sigma$ (DUNE~2028).
\end{itemize}

\begin{keybox}
Die nächsten Jahre werden zeigen, ob bei $5\sigma$ ein neues Verständnis
der Raumzeit-Hardware beginnt.
\end{keybox}

% ── Literatur ────────────────────────────────────────────────
\newpage
\section*{Literaturverzeichnis}
\addcontentsline{toc}{section}{Literaturverzeichnis}

\begin{thebibliography}{99}

\bibitem{NTT2025}
M.\,Stein et al.\ (NTT/Cornell/Stanford),
\textit{Programmable photonics advanced with lithium niobate based waveguide},
Yole Group Industry News (Dezember 2025).

\bibitem{Assumpcao2024}
D.\,Assumpcao et al.,
\textit{A thin film lithium niobate near-infrared platform for multiplexing quantum nodes},
Nature Communications, Vol.\,15 (2024).
\href{https://doi.org/10.1038/s41467-024-54541-2}{doi:10.1038/s41467-024-54541-2}

\bibitem{CERN2026a}
LHCb Collaboration,
\textit{Comprehensive analysis of local and nonlocal amplitudes in the
$B^0 \to K^{*0}\mu^+\mu^-$ decay},
Physical Review Letters (2026),
DOI: 10.1103/24g9-yn9d,
arXiv:2512.18053.

\bibitem{CERN2026b}
LHCb/CMS Collaboration,
\textit{Combined analysis of rare B-meson decays},
LHCb Presse (April 2026).

\bibitem{Arslan2026}
M.\,Arslan et al.,
\textit{The Angular Observables of $\Lambda_b \to \Lambda_c(\to \Lambda^0\pi^+)
\tau^-(\to\pi^-\nu_\tau)\bar{\nu}_\tau$ within the Paradigm of FCCC Anomalies},
arXiv:2604.24922 (April 2026).

\bibitem{Aynbund2025}
A.\,Aynbund, V.\,V.\,Kiselev,
\textit{Effective Field Theory Links Gauge Coupling to Sub-Planckian Cut-off Scale of GeV},
Quantum Zeitgeist (August 2025).

\bibitem{subPlanck2025}
L.\,M.\,de Souza,
\textit{Probing Sub-Planckian Scales via OLQEM and Fibonacci Quasiperiodicity:
A Rigorous Mathematical Verification},
Preprint (2025). [Nicht peer-reviewed; konzeptueller Bezug.]

\bibitem{RGB2025}
A.\,Shimura et al.\ (TDK Corporation),
\textit{Visible light red, green, and blue multiplexer by sputter-deposited
thin-film lithium niobate},
Advanced Photonics Nexus, Vol.\,4, Iss.\,5, 056001 (2025).
\href{https://doi.org/10.1117/1.APN.4.5.056001}{doi:10.1117/1.APN.4.5.056001}

\end{thebibliography}
